% =============================================================================
% Rapport de Projet : Ray Tracer CUDA/CPU - Version condensée HPC
% =============================================================================
\documentclass[12pt,a4paper]{article}

% --- Packages ---
\usepackage[utf8]{inputenc}
\usepackage[T1]{fontenc}
\usepackage[french]{babel}
\usepackage{geometry}
\usepackage{graphicx}
\usepackage{listings}
\usepackage{xcolor}
\usepackage{amsmath}
\usepackage{hyperref}
\usepackage{booktabs}
\usepackage{caption}
\usepackage{subcaption}
\usepackage{float}
\usepackage{fancyhdr}
\usepackage{titlesec}
\usepackage{tikz}
\usetikzlibrary{shapes.geometric,arrows.meta,positioning,calc,fit,backgrounds}

% --- Geometry ---
\geometry{margin=2.5cm}

% --- Code style ---
\definecolor{codebg}{RGB}{245,245,248}
\definecolor{codecomment}{RGB}{100,120,140}
\definecolor{codekeyword}{RGB}{0,100,180}
\definecolor{codestring}{RGB}{180,60,40}

\lstdefinestyle{cudastyle}{
    backgroundcolor=\color{codebg},
    basicstyle=\ttfamily\footnotesize,
    keywordstyle=\color{codekeyword}\bfseries,
    commentstyle=\color{codecomment}\itshape,
    stringstyle=\color{codestring},
    breaklines=true,
    frame=single,
    rulecolor=\color{gray!40},
    numbers=left,
    numberstyle=\tiny\color{gray},
    numbersep=8pt,
    tabsize=4,
    showstringspaces=false,
    language=C++,
    morekeywords={__global__,__device__,__host__,__forceinline__,__constant__,__shared__,
                  dim3,curandState,Color,Ray,BVH,HitRecord,
                  Interval,RenderConfig,Camera,Material,
                  MaterialType,CPURandom,Vec3,Point3,AABB,BVHNode},
    literate={->}{$\rightarrow$}1
}
\lstset{style=cudastyle}

% --- Headers ---
\pagestyle{fancy}
\fancyhf{}
\fancyhead[L]{\small Ray Tracer CUDA/CPU}
\fancyhead[R]{\small\thepage}
\renewcommand{\headrulewidth}{0.4pt}

% --- Title ---
\title{
    \vspace{-1cm}
    {\normalsize Université de Reims Champagne-Ardenne --- Master Informatique}\\[0.1cm]
    {\normalsize Module : Programmation GPU}\\[0.4cm]
    \textbf{Ray Tracer Haute Performance}\\[0.3cm]
    \large Architecture CPU/GPU avec CUDA et OpenMP\\[0.2cm]
    \normalsize Projet HPC -- Romeo2025
}
\author{Nicolas Gotman}
\date{Mars 2026}

% =============================================================================
\begin{document}
\maketitle
\thispagestyle{empty}

\vspace{0.5cm}
\begin{abstract}
Ce rapport détaille le développement d'un moteur de ray tracing exploitant le GPU (CUDA) et le CPU (OpenMP). L'accent est mis sur les optimisations CUDA : exploitation de la hiérarchie mémoire (constante, partagée, globale coalescée), réduction parallèle avec primitives warp-level, CUDA streams pour le recouvrement calcul/transfert, et analyse d'occupancy. Le moteur atteint \textbf{1\,479 MRays/s} sur NVIDIA GH200 H100, avec un speedup de \textbf{21--39$\times$} par rapport au CPU ARM Neoverse V2.
\end{abstract}

\tableofcontents
\vspace{0.5cm}

% =============================================================================
\section{Introduction}
% =============================================================================

La synthèse d'images par path tracing offre un réalisme inégalé mais son coût calculatoire la rend difficilement applicable sans accélération matérielle. Pour chaque pixel, il faut lancer des dizaines de rayons (méthode de Monte Carlo), les suivre de rebond en rebond, et accumuler la contribution lumineuse. Ce problème est intrinsèquement parallélisable : chaque pixel est indépendant, ce qui en fait un candidat idéal pour le GPU.

J'ai développé deux chemins de rendu partageant le même algorithme :
\begin{itemize}
    \item \textbf{GPU (CUDA)} : chaque thread calcule un pixel complet, exploitant les milliers de cœurs du GPU.
    \item \textbf{CPU (OpenMP)} : les lignes de l'image sont distribuées sur les cœurs CPU, servant de référence et de version portable.
\end{itemize}

Le moteur a été testé sur mon PC portable (RTX 4070 Laptop) et le cluster HPC Romeo2025 (NVIDIA GH200 H100). Le code source complet est disponible sur : \url{https://github.com/Gotman08/RayTracing.git}

% =============================================================================
\section{Architecture et algorithmes}
% =============================================================================

\subsection{Organisation modulaire}

Le projet est structuré en bibliothèques d'en-têtes (\texttt{.cuh}) organisées par domaine : fondations (Vec3, Ray, AABB), géométrie (Sphere, Plan), matériaux (Lambertien, Métal, Diélectrique), accélération (BVH), caméra et moteurs de rendu GPU/CPU.

\subsection{Algorithme de rendu}

Le cœur du moteur est la fonction \texttt{ray\_color}, qui suit les rebonds d'un rayon de manière \textbf{itérative} (pas de récursion, car la pile d'appels GPU est limitée). À chaque rebond, l'atténuation est multipliée par la couleur du matériau ; si le rayon atteint le ciel, sa couleur est pondérée par l'atténuation accumulée.

Le dispatch des matériaux utilise un \texttt{switch} sur une énumération au lieu de fonctions virtuelles, car les vtables ne sont pas supportées sur GPU. Trois comportements physiques sont implémentés : diffusion lambertienne, réflexion métallique avec rugosité, et réfraction diélectrique (loi de Snell + approximation de Schlick).

\subsection{Structure d'accélération BVH}

Le \textbf{BVH} (\textit{Bounding Volume Hierarchy}) réduit la complexité d'intersection de $O(n)$ à $O(\log n)$ par rayon. Il est construit sur CPU par approche \textit{top-down} (médiane sur l'axe le plus long) et stocké sous forme de \textbf{tableau plat} avec indices entiers (pas de pointeurs), permettant un transfert direct vers le GPU. La traversée est itérative avec une pile locale de 64 éléments.

Les couleurs HDR sont compressées via le \textbf{tone mapping de Reinhard} ($x \mapsto x/(1+x)$) puis converties en sRGB par correction gamma ($\gamma = 2.2$).

% =============================================================================
\section{Implémentation GPU (CUDA)}
% =============================================================================

\subsection{Architecture du kernel de rendu}

Le kernel principal associe \textbf{un thread CUDA à un pixel}. Les threads sont organisés en blocs 2D de $16 \times 16 = 256$ threads :

\begin{lstlisting}[caption={Lancement du kernel de rendu}]
constexpr int BLOCK_SIZE = 16;
dim3 blocks((width + BLOCK_SIZE - 1) / BLOCK_SIZE,
            (height + BLOCK_SIZE - 1) / BLOCK_SIZE);
dim3 threads(BLOCK_SIZE, BLOCK_SIZE);
render_kernel<<<blocks, threads>>>(d_frame_buffer,
    camera, bvh, config, d_rand_states);
\end{lstlisting}

Chaque thread effectue la totalité du travail pour son pixel : génération des rayons, multi-échantillonnage (SPP), accumulation et tone mapping. Cette stratégie maximise l'indépendance et élimine toute synchronisation inter-threads.

\subsection{Gestion de la mémoire GPU}

Le transfert CPU $\rightarrow$ GPU suit un schéma précis : allocation (\texttt{cudaMalloc}), upload des matériaux (H2D), \textbf{fixup des pointeurs}, upload du BVH (H2D), puis download du framebuffer (D2H) après le rendu.

Le point critique est le \textbf{fixup des pointeurs} : chaque primitive contient un \texttt{Material*} host, invalide sur GPU. Mon code parcourt la scène pour substituer \texttt{\&h\_materials[m]} par \texttt{d\_materials + m}. Cette opération ($<$0.01 ms) doit intervenir \textit{après} l'upload des matériaux et \textit{avant} le BVH.

\subsection{Génération de nombres aléatoires}

Le profiling a révélé que \texttt{curand\_init} représentait $\sim$10\% du temps total. Je l'ai remplacé par le \textbf{hachage de Wang}, une initialisation en $O(1)$ :

\begin{lstlisting}[caption={Initialisation rapide via Wang hash}]
__device__ __forceinline__
unsigned int wang_hash(unsigned int seed) {
    seed = (seed ^ 61) ^ (seed >> 16);
    seed *= 9;
    seed = seed ^ (seed >> 4);
    seed *= 0x27d4eb2d;
    seed = seed ^ (seed >> 15);
    return seed;
}
\end{lstlisting}

% =============================================================================
\section{Hiérarchie mémoire CUDA}
% =============================================================================

L'exploitation de la hiérarchie mémoire est cruciale pour les performances GPU. Le tableau~\ref{tab:mem_hierarchy} résume le compromis latence/capacité avec les données du projet.

\begin{table}[H]
\centering
\small
\caption{Hiérarchie mémoire CUDA et utilisation dans le projet}
\begin{tabular}{llll}
\toprule
\textbf{Niveau} & \textbf{Latence} & \textbf{Portée} & \textbf{Données du projet} \\
\midrule
Registres & $\sim$0 cycles & Thread & ray, rec, attenuation, stack[64] \\
\texttt{\_\_shared\_\_} & $\sim$5 cycles & Bloc & sdata[256] (réduction luminance) \\
\texttt{\_\_constant\_\_} & $\sim$5 cycles & Grille & Camera (120 B), RenderConfig (64 B) \\
L1 / Texture & $\sim$5 cycles & SM & cache automatique (accès BVH) \\
Globale (VRAM) & 400+ cycles & Device & framebuffer, BVH, curandState \\
\bottomrule
\end{tabular}
\label{tab:mem_hierarchy}
\end{table}

\subsection{Mémoire constante : Camera et RenderConfig}

La \texttt{Camera} ($\sim$120 octets) et la \texttt{RenderConfig} ($\sim$64 octets) sont lues \textbf{identiquement par tous les threads}. C'est le cas d'usage idéal de la mémoire constante : le cache constant diffuse la valeur à tout le warp en un seul cycle (broadcast), contre 400--500 cycles en mémoire globale.

\begin{lstlisting}[caption={Utilisation de la mémoire constante}]
__constant__ Camera d_const_camera;
__constant__ RenderConfig d_const_config;
// Copie avant le lancement du kernel
cudaMemcpyToSymbol(d_const_camera, &camera,
                   sizeof(Camera));
\end{lstlisting}

Taille totale : $\sim$184 octets, bien en dessous de la limite de 64 KB.

\subsection{Mémoire partagée : réduction de luminance}

La mémoire \texttt{\_\_shared\_\_} est utilisée dans le kernel de réduction parallèle (section~\ref{sec:reduction}) : chaque bloc de 256 threads utilise un tableau partagé de 1~KB ($256 \times 4$ octets) pour effectuer une réduction en arbre binaire avec \texttt{\_\_syncthreads()}.

\subsection{Coalescence des accès mémoire}

La \textbf{coalescence} permet de fusionner les accès mémoire d'un warp en une seule transaction de 128 octets. L'index linéaire du pixel est $j \times W + i$, avec les threads organisés selon $x$ (dimension rapide du bloc $16\times16$). Les threads $t_0, \ldots, t_{15}$ d'un warp correspondent à des colonnes consécutives sur la même ligne : leurs adresses dans le framebuffer et le tableau \texttt{curandState} sont \textbf{contiguës}, garantissant une bande passante maximale. Un accès par colonnes ($j$ comme dimension rapide) produirait un stride $W$, réduisant le débit d'un facteur pouvant dépasser $\times$10.

% =============================================================================
\section{Réduction parallèle}
\label{sec:reduction}
% =============================================================================

Pour calculer la luminance moyenne (nécessaire au tone mapping adaptatif), j'implémente une \textbf{réduction parallèle} en trois étapes : chaque thread charge la luminance de son pixel en mémoire partagée, réduction intra-bloc par divisions successives avec \texttt{\_\_syncthreads()}, puis le thread 0 accumule via \texttt{atomicAdd}.

\begin{lstlisting}[caption={Kernel de réduction parallèle}]
__global__ void compute_avg_luminance(
    const Color* frame_buffer,
    float* d_total_luminance, int num_pixels) {
    __shared__ float sdata[256];
    int tid = threadIdx.x;
    int gid = blockIdx.x * blockDim.x + threadIdx.x;
    sdata[tid] = (gid < num_pixels) ?
        luminance(frame_buffer[gid]) : 0.0f;
    __syncthreads();
    for (unsigned int s = blockDim.x/2; s > 0; s >>= 1) {
        if (tid < s) sdata[tid] += sdata[tid + s];
        __syncthreads();
    }
    if (tid == 0) atomicAdd(d_total_luminance, sdata[0]);
}
\end{lstlisting}

\subsection{Optimisation warp-level : \texttt{\_\_shfl\_down\_sync}}

Les 5 dernières itérations de la réduction (strides $\leq 32$) n'impliquent qu'un seul warp. Dans un warp, les threads s'exécutent en lockstep : \texttt{\_\_syncthreads()} y est redondant. L'instruction \texttt{\_\_shfl\_down\_sync} échange des valeurs de \textbf{registres} directement entre threads du même warp, sans passer par la mémoire partagée :

\begin{lstlisting}[caption={Réduction warp-level avec \_\_shfl\_down\_sync}]
if (tid < 32) {
    float val = sdata[tid] + sdata[tid + 32];
    val += __shfl_down_sync(0xffffffff, val, 16);
    val += __shfl_down_sync(0xffffffff, val, 8);
    val += __shfl_down_sync(0xffffffff, val, 4);
    val += __shfl_down_sync(0xffffffff, val, 2);
    val += __shfl_down_sync(0xffffffff, val, 1);
    if (tid == 0) atomicAdd(d_total_luminance, val);
}
\end{lstlisting}

Le masque \texttt{0xffffffff} indique que les 32 threads participent. Cette optimisation supprime 5 barrières \texttt{\_\_syncthreads()} et 5 séries d'accès mémoire partagée.

\subsection{Comparaison avec les librairies CUDA}

La librairie \textbf{Thrust} propose \texttt{thrust::transform\_reduce} qui effectue la même réduction en une ligne. \textbf{Avantage du kernel custom} : contrôle fin sur la shared memory et les registres, optimisations spécifiques (\texttt{\_\_shfl\_down\_sync}). \textbf{Avantage de Thrust} : code plus concis, optimisé par NVIDIA, portable entre architectures.

Concernant les autres librairies du toolkit CUDA : \textbf{cuRAND} est utilisée pour la génération de nombres aléatoires (\texttt{curandState} par pixel), bien que l'initialisation ait été remplacée par le Wang hash (section~3.3). \textbf{cuBLAS} (algèbre linéaire dense) n'est pas directement applicable ici car les opérations sont \textit{per-ray} (pas de grandes matrices), mais pourrait servir pour des transformations géométriques par batch d'objets instanciés.

% =============================================================================
\section{CUDA Streams et divergence de warp}
% =============================================================================

\subsection{Recouvrement calcul/transfert avec 2 streams}

Deux \texttt{cudaStream\_t} sont utilisés pour recouvrir les transferts mémoire avec le calcul :
\begin{itemize}
    \item \texttt{stream\_compute} : kernels (init RNG, rendu, réduction, tone mapping)
    \item \texttt{stream\_transfer} : transferts mémoire (H2D matériaux/BVH, D2H framebuffer)
\end{itemize}

Le \textbf{premier recouvrement} superpose l'upload des matériaux (\texttt{cudaMemcpyAsync} sur \texttt{stream\_transfer}) avec l'initialisation curand (\texttt{stream\_compute}). Le \textbf{second} lance le téléchargement asynchrone du framebuffer (D2H) pendant que le CPU prépare la sauvegarde. Le gain est modeste car le kernel domine ($\sim$84\%), mais la technique démontre la maîtrise de l'exécution asynchrone. Une synchronisation (\texttt{cudaStreamSynchronize}) sépare les deux phases.

\subsection{Divergence de warp}

Sur GPU NVIDIA, les 32 threads d'un warp s'exécutent en lockstep. Lorsqu'ils prennent des branches différentes, les deux chemins sont sérialisés. Trois sources de divergence ont été identifiées :

\begin{enumerate}
    \item \textbf{Dispatch de matériaux} : le \texttt{switch(mat.type)} cause de la divergence quand des pixels voisins touchent des matériaux différents. L'impact est maximal aux frontières entre objets.
    \item \textbf{Traversée BVH} : le test \texttt{if (node.is\_leaf)} diverge car les rayons d'un même warp atteignent des profondeurs différentes dans l'arbre.
    \item \textbf{Guard de bordure} : \texttt{if (i >= width || j >= height)} n'affecte que les threads de bord (impact négligeable).
\end{enumerate}

L'utilisation d'un dispatch par \texttt{enum/switch} au lieu de fonctions virtuelles permet au compilateur de mieux prédire les branches. Un tri des rayons par matériau pourrait réduire la divergence du cas 1, mais au prix d'une synchronisation globale non justifiée pour notre scène.

% =============================================================================
\section{Implémentation CPU (OpenMP)}
% =============================================================================

La version CPU partage le même algorithme que le GPU. La parallélisation distribue les lignes de l'image :

\begin{lstlisting}[caption={Parallélisation OpenMP}]
#pragma omp parallel for schedule(dynamic, 1)
for (int j = 0; j < height; j++) {
    CPURandom rng(j * 1337 + 42);
    for (int i = 0; i < width; i++) {
        // ... meme algorithme que le GPU
    }
}
\end{lstlisting}

L'ordonnancement \texttt{schedule(dynamic, 1)} assure un bon équilibrage de charge. Chaque thread utilise son propre générateur \texttt{std::mt19937}. Cette version sert à la fois de version portable et de \textit{ground truth} pour valider le rendu GPU.

% =============================================================================
\section{Analyse de performance et speedup}
% =============================================================================

\subsection{Protocole de mesure}

Les métriques sont collectées via le flag \texttt{--profile} :
\begin{itemize}
    \item \textbf{MRays/s} : $\frac{W \times H \times SPP}{temps}$ (rayons primaires par seconde)
    \item \textbf{Speedup} : $S = T_{CPU} / T_{GPU}$
\end{itemize}

\textbf{Plateformes de test} : (1) PC portable --- RTX 4070 Laptop (5888 cœurs CUDA, 8 Go GDDR6X), CPU x86 20 cœurs via WSL2 ; (2) Cluster Romeo --- GH200 H100 (120 Go HBM3, architecture Hopper), CPU ARM Neoverse V2 (16 threads OpenMP), CUDA 12.6.

\subsection{Résultats}

Sur le \textbf{PC portable}, le throughput GPU atteint $\sim$270 MRays/s en 1080p, avec un speedup de \textbf{73--165$\times$} (CPU limité par WSL2). Sur le \textbf{GH200}, le throughput culmine à \textbf{1\,479 MRays/s} en 4K ($\textbf{5.3}\times$ plus rapide que la RTX 4070), avec un speedup GPU/CPU de \textbf{21--39$\times$} (tableaux~\ref{tab:benchmarks_romeo_gpu} et \ref{tab:benchmarks_romeo_cpu} en annexe).

\begin{figure}[H]
    \centering
    \includegraphics[width=0.85\textwidth]{../output/chart_gpu_comparison.png}
    \caption{Comparaison des performances GPU : RTX 4070 Laptop vs GH200 H100 (MRays/s)}
    \label{fig:gpu_comparison}
\end{figure}

\begin{figure}[H]
    \centering
    \includegraphics[width=0.85\textwidth]{../output/chart_cpu_vs_gpu.png}
    \caption{Temps de rendu CPU vs GPU sur GH200 (échelle logarithmique)}
    \label{fig:cpu_vs_gpu}
\end{figure}

\subsection{Profiling GPU}

\begin{table}[H]
\centering
\caption{Décomposition du temps GPU (800$\times$600, 100 SPP)}
\begin{tabular}{lcc}
\toprule
\textbf{Étape} & \textbf{Temps (ms)} & \textbf{Proportion} \\
\midrule
Allocation mémoire GPU & 3.56 & 1.37\% \\
Upload matériaux (H2D) & 1.78 & 0.69\% \\
Construction BVH (CPU) & 0.15 & 0.06\% \\
Upload BVH (H2D) & 0.19 & 0.07\% \\
Init états aléatoires & 24.13 & 9.31\% \\
\textbf{Kernel de rendu} & \textbf{219.06} & \textbf{84.49\%} \\
Download framebuffer (D2H) & 10.31 & 3.98\% \\
\bottomrule
\end{tabular}
\label{tab:profiling}
\end{table}

Le kernel domine ($\sim$84\%), confirmant que les transferts ne sont pas un goulot d'étranglement. Sur le GH200 en 1080p @ 500 SPP, les transferts H2D + D2H restent $<$0.2\%, cohérent avec la bande passante HBM3 (4 To/s). Le profilage interne a été préféré à \textbf{Nsight Compute} (successeur de \texttt{nvprof}) en raison des limitations sous WSL2, mais \texttt{cudaGetLastError()} et \textbf{Compute Sanitizer} (\texttt{compute-sanitizer --tool memcheck}) ont été utilisés systématiquement pour détecter les accès mémoire invalides.

\subsection{Occupancy et caractérisation compute-bound}

L'occupancy (ratio warps actifs / warps max par SM) est calculée via l'API CUDA :

\begin{lstlisting}[caption={Calcul d'occupancy}]
int num_blocks_per_sm;
cudaOccupancyMaxActiveBlocksPerMultiprocessor(
    &num_blocks_per_sm, render_kernel,
    BLOCK_SIZE * BLOCK_SIZE, 0);
\end{lstlisting}

La bande passante effective est estimée à $\sim$64 octets/pixel (lecture \texttt{curandState} + écriture \texttt{Color}). Elle reste bien en dessous du pic théorique (272 Go/s pour RTX 4070, 4 To/s pour GH200), confirmant que le kernel est \textbf{compute-bound}. Le throughput sature à $\sim$1\,450 MRays/s au-delà de 1080p : en basse résolution, les 132 SM du H100 ne sont pas entièrement saturés ; à partir de 1080p, le GPU est pleinement exploité.

\subsection{Discussion : sensibilité du speedup à la baseline}

Le speedup varie fortement : 73--165$\times$ sur le laptop contre 21--39$\times$ sur Romeo. Le GH200 est pourtant 5.3$\times$ plus rapide en absolu --- la différence vient du CPU de référence : le Neoverse V2 atteint 33--36 MRays/s contre 1.7--3.0 MRays/s sous WSL2 (surcoût Hyper-V, caches serveur vs. mobile, DDR5 native). \textbf{Un speedup GPU doit toujours être interprété en regard de la baseline CPU choisie.}

% =============================================================================
\section{Conclusion}
% =============================================================================

Ce projet a confronté les réalités de la programmation GPU haute performance à travers un path tracer Monte Carlo, atteignant un speedup de 21--39$\times$ sur le GH200 et un throughput de 1\,479 MRays/s en 4K.

Les concepts du cours mis en pratique couvrent : l'exploitation de la \textbf{hiérarchie mémoire} (constante pour le broadcast warp, partagée pour la réduction, globale coalescée), la \textbf{réduction parallèle} avec primitives warp-level (\texttt{\_\_shfl\_down\_sync}), les \textbf{CUDA Streams} pour le recouvrement calcul/transfert, l'analyse de \textbf{divergence de warp}, la caractérisation \textbf{compute-bound} via occupancy et bande passante, et la comparaison avec les \textbf{librairies CUDA} (Thrust, cuRAND, cuBLAS). Le débogage a reposé sur \texttt{cudaGetLastError()} et Compute Sanitizer.

En perspective, l'importance sampling et la construction BVH sur GPU (LBVH) permettraient d'améliorer convergence et performances.

% =============================================================================
\appendix
\section{Données brutes de benchmark}
% =============================================================================

\begin{table}[H]
\centering
\caption{Performances GPU sur NVIDIA GH200 H100 (Romeo)}
\begin{tabular}{lcccc}
\toprule
\textbf{Résolution} & \textbf{SPP} & \textbf{Temps (s)} & \textbf{MRays/s} \\
\midrule
640$\times$480 & 100 & 0.045 & 683 \\
1280$\times$720 & 100 & 0.073 & 1\,262 \\
1920$\times$1080 & 100 & 0.147 & 1\,411 \\
1920$\times$1080 & 500 & 0.719 & 1\,442 \\
3840$\times$2160 & 100 & 0.569 & 1\,458 \\
3840$\times$2160 & 500 & 2.804 & 1\,479 \\
\bottomrule
\end{tabular}
\label{tab:benchmarks_romeo_gpu}
\end{table}

\begin{table}[H]
\centering
\caption{Performances CPU sur GH200 ARM Neoverse V2 (16 threads)}
\begin{tabular}{lccccc}
\toprule
\textbf{Résolution} & \textbf{SPP} & \textbf{CPU (s)} & \textbf{CPU MRays/s} & \textbf{GPU (s)} & \textbf{Speedup} \\
\midrule
640$\times$480 & 100 & 0.933 & 33 & 0.045 & 21$\times$ \\
1280$\times$720 & 100 & 2.558 & 36 & 0.073 & 35$\times$ \\
1920$\times$1080 & 100 & 5.805 & 36 & 0.147 & 39$\times$ \\
\bottomrule
\end{tabular}
\label{tab:benchmarks_romeo_cpu}
\end{table}

\end{document}
